\section*{Aufgabenstellung}
Ermitteln Sie mit Binwalk den Inhalt eines squashfs Containers am Ende des Server Images.

\section*{Lösung}

Im ersten Schritt wurden die letzten 50MB des Images extrahiert, da hier der squashfs Container vermutet wird:
Auf diesem extrahierten Image wurde dann mit Binwalk der Inhalt des squashfs Containers extrahiert:
\begin{lstlisting}[language=bash]
root@Alpha-17:/home/luis/Programming/GitHub/University/DigitalForensic# dd if=/mnt/ewf_server/ewf1 of=/tmp/tail.bin bs=1M skip=$(($(stat -c%s /mnt/ewf_server/ewf1)/1048576 - 50))
50+1 records in
50+1 records out
53243904 bytes (53 MB, 51 MiB) copied, 0.0928973 s, 573 MB/s
root@Alpha-17:/home/luis/Programming/GitHub/University/DigitalForensic# binwalk /tmp/tail.bin

DECIMAL       HEXADECIMAL     DESCRIPTION
--------------------------------------------------------------------------------
49283072      0x2F00000       Squashfs filesystem, little endian, version 4.0, compression:gzip, size: 3960790 bytes, 3722 inodes, blocksize: 131072 bytes, created: 2026-03-23 12:36:31

root@Alpha-17:/home/luis/Programming/GitHub/University/DigitalForensic# dd if=/tmp/tail.bin of=/tmp/squashfs.img bs=1 skip=49283072
3960832+0 records in
3960832+0 records out
3960832 bytes (4.0 MB, 3.8 MiB) copied, 5.8737 s, 674 kB/s
root@Alpha-17:/home/luis/Programming/GitHub/University/DigitalForensic# dd if=/tmp/tail.bin of=/tmp/squashfs.img bs=1 skip=49283072
3960832+0 records in
3960832+0 records out
3960832 bytes (4.0 MB, 3.8 MiB) copied, 5.85739 s, 676 kB/s
root@Alpha-17:/home/luis/Programming/GitHub/University/DigitalForensic# file /tmp/squashfs.img
/tmp/squashfs.img: Squashfs filesystem, little endian, version 4.0, zlib compressed, 3960790 bytes, 3722 inodes, blocksize: 131072 bytes, created: Mon Mar 23 12:36:31 2026
\end{lstlisting}

Auszüge aus dem extrahierten squashfs Container:
\begin{lstlisting}[language=bash]
root@Alpha-17:/home/luis/Programming/GitHub/University/DigitalForensic# unsquashfs -l /tmp/squashfs.img
squashfs-root/.editorconfig
squashfs-root/.git
squashfs-root/4n622tjddwejdqmb4qsd5zyjpwh7f6awnvdhostckei5tui6ohdla3id.onion
squashfs-root/4n622tjddwejdqmb4qsd5zyjpwh7f6awnvdhostckei5tui6ohdla3id.onion/hostname
squashfs-root/4n622tjddwejdqmb4qsd5zyjpwh7f6awnvdhostckei5tui6ohdla3id.onion/hs_ed25519_public_key
squashfs-root/4n622tjddwejdqmb4qsd5zyjpwh7f6awnvdhostckei5tui6ohdla3id.onion/hs_ed25519_secret_key
\end{lstlisting}

Es handelt sich definitiv um ein git Repository, welches die Konfiguration eines Tor Hidden Services enthält.
Es gibt 3316 Dateien mit der Endung .onion:

\begin{lstlisting}[language=bash]
root@Alpha-17:/home/luis/Programming/GitHub/University/DigitalForensic# unsquashfs -l /tmp/squashfs.img | grep -i '.onion' | wc -l
3316
\end{lstlisting}

Es könnte sich um das Git Repository mkp224o handeln.
Ein Tool das zur generierung von Tor Hidden Services genutzt wird.
